\RequirePackage[l2tabu, orthodox]{nag}
\RequirePackage{silence}
\WarningFilter{nag}{There is no environment ``centering'' }%nag complains because beamer titlepage uses a centering environment
\WarningFilter{nag}{1 complaints in total}
\WarningFilter{ifpdf}{Someone has redefined \pdfoutput}
\WarningFilter{fmtcount}{\ordinal already defined use \FCordinal instead}
\documentclass[english, french]{beamer}
%INSTALL

%avoids a warning
\usepackage[log-declarations=false]{xparse}
\usepackage{fontspec} %font selecting commands
\usepackage{xunicode}
%warn about missing characters
\tracinglostchars=2

%REDAC
\usepackage{booktabs}
\usepackage{calc}

\usepackage{mathtools} %load this before babel!
	\mathtoolsset{showonlyrefs,showmanualtags}

\usepackage{babel}
%suppresses the warning about frenchb not modifying the captions (“—” to “:” in “Figure 1 – Legend”).
	\frenchbsetup{AutoSpacePunctuation=false,SuppressWarning=false}

\usepackage[super]{nth}
\usepackage{listings} %typeset source code listings
	\lstset{language=XML,tabsize=2,literate={"}{{\tt"}}1,captionpos=b}
\usepackage[nolist,smaller,printonlyused]{acronym}%,smaller option produces warnings from relsize in some cases, it seems.
\usepackage[nodayofweek]{datetime}%must be loaded after the babel package
\usepackage{xspace}
\usepackage{hyperref}% option pdfusetitle must be introduced here, not in hypersetup.
%breaklinks makes links on multiple lines into different PDF links to the same target.
%colorlinks (false): Colors the text of links and anchors. The colors chosen depend on the the type of link. In spite of colored boxes, the colored text remains when printing.
%linkcolor=black: this leaves other links in colors, e.g. refs in green, don't print well.
%pdfborder (0 0 1, set to 0 0 0 if colorlinks): width of PDF link border
%hidelinks
\hypersetup{breaklinks,bookmarksopen,colorlinks=true,urlcolor=blue,linkcolor=,hyperfigures=true}
% hyperref doc says: Package bookmark replaces hyperref’s bookmark organization by a new algorithm (...) Therefore I recommend using this package.
\usepackage{bookmark}

% center floats by default, but do not use with float
% \usepackage{floatrow}
% \makeatletter
% \g@addto@macro\@floatboxreset\centering
% \makeatother
\usepackage{ragged2e} %new com­mands \Cen­ter­ing, \RaggedLeft, and \RaggedRight and new en­vi­ron­ments Cen­ter, FlushLeft, and FlushRight, which set ragged text and are eas­ily con­fig­urable to al­low hy­phen­ation (the cor­re­spond­ing com­mands in LaTeX, all of whose names are lower-case, pre­vent hy­phen­ation al­to­gether). 
\usepackage{siunitx} %[expproduct=tighttimes, decimalsymbol=comma]
\sisetup{detect-all}% to detect e.g. when in math mode (use a math font)
\usepackage{braket} %for \Set
\usepackage{natbib}

\usepackage{amsmath,amsthm}
% \usepackage{amsfonts} %not required?
% \usepackage{dsfont} %for what?
%unicode-math overwrites the following commands from the mathtools package: \dblcolon, \coloneqq, \Coloneqq, \eqqcolon. Using the other colon-like commands from mathtools will lead to inconsistencies. Plus, Using \overbracket and \underbracke from mathtools package. Use \Uoverbracket and \Uunderbracke for original unicode-math definition.
%use exclusively \mathbf and choose math bold style below.
\usepackage[warnings-off={mathtools-colon, mathtools-overbracket}, bold-style=ISO]{unicode-math}

\defaultfontfeatures{
	Fractions=On,
	Mapping=tex-text% to turn "--" into dashes, useful for bibtex%%
}
\defaultfontfeatures[\rmfamily, \sffamily]{
	Fractions=On,
	Mapping=% to leave " alone (disable the default mapping tex-text; requires loading the font afterwards?)
}
\newfontfamily\xitsfamily{XITS}
\newfontfamily\texgyretermesfamily{TeX Gyre Termes}
\newfontfamily\texgyreherosfamily{TeX Gyre Heros}
\newfontfamily\lmfamily{Latin Modern Roman}
\setmainfont{Latin Modern Roman}
\setsansfont{Latin Modern Sans}
\setmonofont{Latin Modern Mono}
\defaultfontfeatures{}%disable default font features to avoid warnings with math fonts.
\setmathfont{XITS Math}
\setmathfont[range={\mathcal,\mathbfcal},StylisticSet=1]{XITS Math}

\usepackage{cleveref}% cleveref should go "laster" than hyperref
%GRAPHICS
\usepackage{pgf}
\usepackage{pgfplots}
	\usetikzlibrary{matrix,fit,plotmarks,calc,trees,shapes.geometric,positioning,plothandlers}
\pgfplotsset{compat=1.11}
\usepackage{graphicx}

\graphicspath{{graphics/},{graphics-dm/}}
\DeclareGraphicsExtensions{.pdf}

%HACKING
\usepackage{printlen}
\uselengthunit{mm}
% 	\newlength{\templ}% or LenTemp?
% 	\setlength{\templ}{6 pt}
% 	\printlength{\templ}
\usepackage{etoolbox} %for addtocmd command
\usepackage{scrhack}% load at end. Corrects a bug in float package, which is outdated but might be used by other packages
\usepackage{xltxtra} %somebody said that this is loaded by fontspec, but does not seem correct: if not loaded explicitly, does not appear in the log and \showhyphens is not corrected.

%Beamer-specific
%ADD
\usepackage{appendixnumberbeamer}
%\setbeamersize{text margin left=0.1cm, text margin right=0.1cm} 
\setbeamertemplate{navigation symbols}{}
\usetheme{BrusselsBelgium}
\usefonttheme{professionalfonts}


\newcommand{\R}{ℝ}
\newcommand{\N}{ℕ}
\newcommand{\Z}{ℤ}
\newcommand{\card}[1]{\lvert{#1}\rvert}
\newcommand{\powerset}[1]{\mathscr{P}(#1)}%\mathscr rather than \mathcal: scr is rounder than cal (at least in XITS Math).
\newcommand{\suchthat}{\;\ifnum\currentgrouptype=16 \middle\fi|\;}
%\newcommand{\Rplus}{\reels^+\xspace}

\AtBeginDocument{%
	\renewcommand{\epsilon}{\varepsilon}
% we want straight form of \phi for mathematics, as recommended in UTR #25: Unicode support for mathematics.
%	\renewcommand{\phi}{\varphi}
}

% with amssymb, but I don’t want to use amssymb just for that.
% \newcommand{\restr}[2]{{#1}_{\restriction #2}}
%\newcommand{\restr}[2]{{#1\upharpoonright}_{#2}}
\newcommand{\restr}[2]{{#1|}_{#2}}%sometimes typed out incorrectly within \set.
%\newcommand{\restr}[2]{{#1}_{\vert #2}}%\vert errors when used within \Set and is typed out incorrectly within \set.
\DeclareMathOperator*{\argmax}{arg\,max}
\DeclareMathOperator*{\argmin}{arg\,min}


%ARG TH
\newcommand{\AF}{\mathcal{AF}}
\newcommand{\labelling}{\mathcal{L}}
\newcommand{\labin}{\textbf{in}\xspace}
\newcommand{\labout}{\textbf{out}}
\newcommand{\labund}{\textbf{undec}\xspace}
\newcommand{\nonemptyor}[2]{\ifthenelse{\equal{#1}{}}{#2}{#1}}
\newcommand{\gextlab}[2][]{
	\labelling{\mathcal{GE}}_{(#2, \nonemptyor{#1}{\ibeatsr{#2}})}
}
\newcommand{\allargs}{A^*}
\newcommand{\args}{A}
\newcommand{\ar}{a}

%MCDA+Arg
\newcommand{\dm}{d}
\newcommand{\ileadsto}{\leadsto}
\newcommand{\mleadsto}[1][\eta]{\leadsto_{#1}}
\newcommand{\ibeats}{\vartriangleright}
\newcommand{\mbeats}[1][\eta]{\vartriangleright_{#1}}


%MISC
\newcommand{\lequiv}{\Vvdash}
\newcommand{\weightst}{W^{\,t}}

%MCDA classical
\newcommand{\crits}{\mathcal{J}}
\newcommand{\altspace}{\mathbb{A}}
\newcommand{\alts}{A}

%Sorting
\newcommand{\cats}{\mathcal{C}}
\newcommand{\catssubsets}{2^\cats}
\newcommand{\catgg}{\vartriangleright}
\newcommand{\catll}{\vartriangleleft}
\newcommand{\catleq}{\trianglelefteq}
\newcommand{\catgeq}{\trianglerighteq}
\newcommand{\alttoc}[2][x]{(#1 \xrightarrow{} #2)}
\newcommand{\alttocat}[3]{(#2 \xrightarrow{#1} #3)}
\newcommand{\alttoI}{(x \xrightarrow{} \left[\underline{C_x}, \overline{C_x}\right])}
\newcommand{\alttocatdm}[3][t]{\left(#2 \thinspace \raisebox{-3pt}{$\xrightarrow{#1}$}\thinspace #3\right)}
\newcommand{\alttocatatleast}[2]{\left(#1 \thinspace \raisebox{-3pt}{$\xrightarrow[]{≥}$}\thinspace #2\right)}
\newcommand{\alttocatatmost}[2]{\left(#1 \thinspace \raisebox{-3pt}{$\xrightarrow[]{≤}$}\thinspace #2\right)}

\newcommand{\source}{\scriptsize}
\newcommand{\commentOC}[1]{{\selectlanguage{french}{\todo{OC : #1}}}}
%Or: \todo[color=green!40]

%this probably requires outdated float package, see doc KomaScript for an alternative.
% \newfloat{program}{t}{lop}
% \floatname{program}{PM}

%\crefname{axiom}{axiom}{axioms}%might be needed for workaround bug in cref when defining new theorems?

%\ifdefined\theorem\else
%\newtheorem{theorem}{\iflanguage{english}{Theorem}{Théorème}}
%\fi

%which line breaks are chosen: accept worse lines, therefore reducing risk of overfull lines. Default = 200
\tolerance=2000
%accept overfull hbox up to...
\hfuzz=2cm
%reduces verbosity about the bad line breaks
\hbadness 5000
%sloppy sets tolerance to 9999
\apptocmd{\sloppy}{\hbadness 10000\relax}{}{}

% WRITING
%\newcommand{\ie}{i.e.\@\xspace}%to try
%\newcommand{\eg}{e.g.\@\xspace}
%\newcommand{\etal}{et al.\@\xspace}
\newcommand{\ie}{i.e.\ }
\newcommand{\eg}{e.g.\ }
\newcommand{\mkkOK}{\checkmark}%\color{green}{\checkmark}
\newcommand{\mkkREQ}{\ding{53}}%requires pifont?%\color{green}{\checkmark}
\newcommand{\mkkNO}{}%\text{\color{red}{\textsf{X}}}

\makeatletter
\newcommand{\boldor}[2]{%
	\ifnum\strcmp{\f@series}{bx}=\z@
		#1%
	\else
		#2%
	\fi
}
\newcommand{\textstyleElProm}[1]{\boldor{\MakeUppercase{#1}}{\textsc{#1}}}
\makeatother
\newcommand{\electre}{\textstyleElProm{Électre}\xspace}
\newcommand{\electreIv}{\textstyleElProm{Électre Iv}\xspace}
\newcommand{\electreIV}{\textstyleElProm{Électre IV}\xspace}
\newcommand{\electreIII}{\textstyleElProm{Électre III}\xspace}
\newcommand{\electreTRI}{\textstyleElProm{Électre Tri}\xspace}
% \newcommand{\utadis}{\texorpdfstring{\textstyleElProm{utadis}\xspace}{UTADIS}}
% \newcommand{\utadisI}{\texorpdfstring{\textstyleElProm{utadis i}\xspace}{UTADIS I}}

%TODO
% \newcommand{\textstyleElProm}[1]{{\rmfamily\textsc{#1}}} 


\newlength{\GraphsNodeSep}
\setlength{\GraphsNodeSep}{7mm}

% MCDA Drawing Sorting
\newlength{\MCDSCatHeight}
\setlength{\MCDSCatHeight}{6mm}
\newlength{\MCDSAltHeight}
\setlength{\MCDSAltHeight}{4mm}
%separation between two vertical alts
\newlength{\MCDSAltSep}
\setlength{\MCDSAltSep}{2mm}
\newlength{\MCDSCatWidth}
\setlength{\MCDSCatWidth}{3cm}
\newlength{\MCDSEvalRowHeight}
\setlength{\MCDSEvalRowHeight}{6mm}
\newlength{\MCDSAltsToCatsSep}
\setlength{\MCDSAltsToCatsSep}{1.5cm}
\newcounter{MCDSNbAlts}
\newcounter{MCDSNbCats}
\newlength{\MCDSArrowDownOffset}
\setlength{\MCDSArrowDownOffset}{0mm}

\tikzset{/Graphs/dot/.style={
	shape=circle, fill=black, inner sep=0, minimum size=1mm
}}
\tikzset{/MC/D/S/alt/.style={
	shape=rectangle, draw=black, inner sep=0, minimum height=\MCDSAltHeight, minimum width=2.5cm, anchor=north east
}}
\tikzset{MC/D/S/pref/.style={
	shape=ellipse, draw=gray, thick
}}
\tikzset{/MC/D/S/cat/.style={
	shape=rectangle, draw=black, inner sep=0, minimum height=\MCDSCatHeight, minimum width=\MCDSCatWidth, anchor=north west
}}
\tikzset{/MC/D/S/evals matrix/.style={
	matrix, row sep=-\pgflinewidth, column sep=-\pgflinewidth, nodes={shape=rectangle, draw=black, inner sep=0mm, text depth=0.5ex, text height=1em, minimum height=\MCDSEvalRowHeight, minimum width=12mm}, nodes in empty cells, matrix of nodes, inner sep=0mm, outer sep=0mm, row 1/.style={nodes={draw=none, minimum height=0em, text height=, inner ysep=1mm}}
}}

% Beliefs
\tikzset{/Beliefs/D/S/attacker/.style={
	shape=rectangle, draw, minimum size=8mm
}}
\tikzset{/Beliefs/D/S/supporter/.style={
	shape=circle, draw
}}

\newcommand{\tikzmark}[1]{%
	\tikz[overlay, remember picture, baseline=(#1.base)] \node (#1) {};%
}


\begin{acronym}
\acro{AMCD}{Aide Multicritère à la Décision}
\acro{ASA}{Argument Strength Assessment}
\acro{DA}{Decision Analysis}
\acro{DM}{Decision Maker}
\acro{DPr}{Deliberated Preferences}
\acro{DRSA}{Dominance-based Rough Set Approach}
\acro{DSS}{Decision Support Systems}
\acrodefplural{DSS}{Decision Support Systems}
% \newacroplural{DSS}[DSSes]{Decision Support Systems}
\acro{EJOR}{European Journal of Operational Research}
\acro{LNCS}{Lecture Notes in Computer Science}
\acro{MCDA}{Multicriteria Decision Aid}
\acro{MIP}{Mixed Integer Program}
\acro{NCSM}{Non Compensatory Sorting Model}
\acro{PL}{Programme Linéaire}
\acro{PLNE}{Programme Linéaire en Nombres Entiers}
\acro{PM}{Programme Mathématique}
\acro{MP}{Mathematical Program}
\acro{MIP}{Mixed Integer Program}
% \newacroplural{PM}{Programmes Mathématiques}
%acrodefplural since version 1.35, my debian has \ProvidesPackage{acronym}[2009/01/25, v1.34, Support for acronyms (Tobias Oetiker)]
\acrodefplural{PM}{Programmes Mathématiques}
\acro{PMML}{Predictive Model Markup Language}
\acro{RESS}{Reliability Engineering \& System Safety}
\acro{SMAA}{Stochastic Multicriteria Acceptability Analysis}
\acro{URPDM}{Uncertainty and Robustness in Planning and Decision Making}
\acro{XML}{Extensible Markup Language}
\end{acronym}

%this approach does not generalize to multipart nodes.
%\tikzset{/uml/interface/.style={rectangle, draw, fill=yellow!20, align=left, node font=\fontspec{Latin Modern Mono Light}, node contents={$<<$interface$>>$\\\textbf{#1}}, name=#1
%}}
%this approach does not apply for <<interface>>. I’ll do it manually, as I can’t find anything better. Use: \nodepart[font=]{one} \small <<interface>>\\\bfseries ItemDAO
\tikzset{/uml/class/.style={rectangle, draw, align=center, fill=yellow!20, node font=\ttfamily, font=\bfseries
}}
\tikzset{/uml/abstract class/.style={rectangle, draw, align=center, fill=yellow!20, node font=\fontspec{Latin Modern Mono Light}, font=\bfseries\itshape
}}
\tikzset{/uml/class3/.style={
rectangle split, rectangle split every empty part={}, rectangle split parts=3, rectangle split part align={center, left, left}, draw, fill=yellow!20, rectangle split empty part height=0, node font=\ttfamily, every one node part/.style={font=\bfseries, align=center}, every two node part/.style={align=left}, every three node part/.style={align=left}
}}
\tikzset{/uml/extends/.style={draw, -open triangle 45}}
\tikzset{/uml/implements/.style={draw, -open triangle 45, dashed}}
%prefix= used to cancel \texttt, otherwise this cancels the effect of the font selection \fontspec{Latin Modern Mono Light}, and thus renders \bfseries inoperant. With prefix=, the result is correct.
%\jeeref[prefix=]{javax.faces.component.html/HtmlCommandButton}\\
\tikzset{/uml/table/.style={rectangle, draw, align=center, fill=blue!20, node font=\ttfamily, font=\bfseries
}}
\tikzset{/uml/table2/.style={
rectangle split, rectangle split every empty part={}, rectangle split parts=2, rectangle split part align={center, left}, draw, fill=blue!20, rectangle split empty part height=0, node font=\ttfamily, every one node part/.style={font=\bfseries, align=center}, every two node part/.style={align=left}
}}
\tikzset{/uml/dbkey/.style={draw, ->}}
%http://tex.stackexchange.com/questions/79781/placing-anchor-before-and-after-text-in-multipart-rectangle



\title{Conception d’applications internet}
\subtitle{JPA}
\subject{object / relationship mismatch}
\keywords{ORM ; transactions ; isolation}
\author{Olivier Cailloux}
\institute[LAMSADE]{LAMSADE, Université Paris-Dauphine}
\date{Version du \today}

\begin{document}
\bibliographystyle{apalike}

\begin{frame}[plain]
	\tikz[remember picture,overlay]{
		\path (current page.south west) node[anchor=south west, inner sep=0] {
			\includegraphics[height=1cm]{LAMSADE95.jpg}
		};
		\path (current page.south) ++ (0, 1mm) node[anchor=south, inner sep=0] {
			\includegraphics[height=9mm]{Dauphine.jpg}
		};
		\path (current page.south east) node[anchor=south east, inner sep=0] {
			\includegraphics[height=1cm]{PSL.png}
		};
	}
   \titlepage
\end{frame}
\addtocounter{framenumber}{-1}

\section{JPA}
\subsection{Introduction}
\begin{frame}
	\frametitle{Java Persistence API : introduction}
	\begin{itemize}
		\item Standard pour gérer la persistence
		\item Pour Java SE et Java EE
		\item Modèle Objet / Relationnel (ORM)
		\item JPA définit les concepts et interfaces, fournisseur JPA les implémente (exemple : Hibernate)
		\item Pour cette présentation : ORM ; Accès concurrents ; JPA
		\item Appuié sur JDBC
	\end{itemize}
\end{frame}

\begin{frame}
	\frametitle{Standards JCP}
	\begin{itemize}
		\item Implication de \og{}la communauté\fg{} pour standards Java
		\item JCP ? \pause Java Community Process \pause
		\item Définit les JSR : standards utilisés en Java SE ou Java EE
		\item Spécifications tiennent compte de nombreux avis d’horizons divers
		\item JSR 338 : JPA 2.1 ; JSR 345 : EJB 3.2 ; JSR 342 : Java EE 7 ; JSR 346 : CDI…
		\item Tentions entre standard ouvert et contrôle ! (2010, Apache \href{https://blogs.apache.org/foundation/entry/the_asf_resigns_from_the}{quitte} le comité JCP ; Doug Lea \href{http://gee.cs.oswego.edu/dl/html/jcp22oct10.html}{également}, en faveur de OpenJDK…)
	\end{itemize}
\end{frame}

\begin{frame}
	\frametitle{Faiblesses du modèle relationnel pur}
	Problème : \og{}Mismatch\fg{} Objet / Relationnel
	\begin{itemize}
		\item Modèle de données (DM) : objets
		\item Références directionnelles
		\item Modèle Entité / Relationnel : sans comportement
		\item Pas de notion de navigation
		\item Chargement automatique ? \pause Éviter de charger tout le graphe !\pause
		\item Problème classique : n+1 select
		\item Cohérence à maintenir entre BD et objets : types, colonnes, contraintes not null ou autres…
		\item Répétition lors écriture des requêtes de base
		\item Difficultés particulières : héritage et autres concepts objet
	\end{itemize}
\end{frame}

\begin{frame}
	\frametitle{Avantages d’une solution ORM}
	ORM ? \pause Object / Relational Mapping \pause
	\begin{itemize}
		\item Détection des modifications avec accès BD optimisés
		\item Réduction code répétitif
		\item Meilleure portabilité
		\item Permet modèle objet fin
		\item Facilite refactoring et développement agile
	\end{itemize}
	{\tiny De : \href{http://www.lamsade.dauphine.fr/~manouvri/HIBERNATE/SLIDES/ORM.pdf}{diapos} Maude Manouvrier, p. 22}
\end{frame}

\subsection{ORM et entités}
\begin{frame}
	\frametitle{Entité}
	\begin{itemize}
		\item Lien entre DM et BD : \emph{entités} et annotations sur entités
		\item DM contient des classes \emph{entités}
		\item Entité : instance représente {\tiny pas toujours} une ligne d’une table
		\item Marquer classe \jeeref[@]{javax.persistence/Entity} {\tiny (voir aussi \jeeref[@]{javax.persistence/Table})}
		\item Marquer champs transiants {\tiny (ou méthodes get*)} \jeeref[@]{javax.persistence/Transient} {\tiny (et persistants \jeeref[@]{javax.persistence/Column} si désiré)}
		\item Marquer un champ persistant id \jeeref[@]{javax.persistence/Id} (et \jeeref[@]{javax.persistence/GeneratedValue})
		\item Id représente la clé primaire
		\item Id initialisé par le fournisseur de persistence (pas de \texttt{setId} public) {\tiny (sauf si clé naturelle, généralement déconseillé)}
		\item Pour permettre concurrence optimiste : marquer un champ persistant version \jeeref[@]{javax.persistence/Version} (écriture : slmt fournisseur)
	\end{itemize}
\end{frame}

\begin{frame}
	\frametitle{Unité de persistance}
	\begin{itemize}
		\item Ensemble d’entités contenu dans une \og{}persistence unit\fg{}
		\item Unité liée à un pilote JDBC et des propriétés de connexion par configuration
		\item Unité associée à \jeeref{javax.persistence/EntityManagerFactory}
		\item \jeeref{javax.persistence/EntityManager} gère entités dans unité persistance
	\end{itemize}
\end{frame}

\subsection{Contexte de persistance}
\begin{frame}
	\frametitle{Contexte de persistance}
	\begin{itemize}
		\item EntityManager lié à un \emph{contexte de persistance}
		\item Contexte : associe une ligne DB à (max.) une instance d’entité
	\end{itemize}
	Cycle typique :
	\begin{itemize}
		\item \texttt{e = new MyEntity(); e.setTruc(…);//} e est \og{}new\fg{}
		\item \texttt{entityManager.persist(e);//} e est \og{}managed\fg{}
		\item fermeture du entityManager : e est \og{}detached\fg{}
	\end{itemize}
	\begin{block}{États d’une instance d’entité}
		\begin{description}[detached]
			\item[new] sans identité persistante, pas dans le contexte
			\item[managed] avec identité persistante, dans le contexte
			\item[detached] avec identité persistante, pas dans le contexte
			\item[removed] avec identité persistante, dans le contexte de persistance, marquée pour effacement
		\end{description}
	\end{block}
\end{frame}

\begin{frame}
	\frametitle{EM et transactions}
	\begin{itemize}
		\item Les changements du contexte de persistance via l’entity manager sont synchronisés avec la BD au commit {\tiny (ou flush)}
		\item Les changements à la BD sont annulés par un rollback
		\item Injection dans EJB : \jeeref[@]{javax.persistence/PersistenceContext} \texttt{EntityManager}
		\item Persistence context (par défaut) a le scope de la transaction en cours
	\end{itemize}
	Exemple de changement :
	\begin{itemize}
		\item \texttt{e=entityManager.find(…);//} e est \og{}managed\fg{}
		\item \texttt{e.setTruc(…);//} changement marqué
		\item fermeture du entityManager : synchronisation
	\end{itemize}
\end{frame}

\subsection{En pratique}
\begin{frame}
	\frametitle{Configuration unité de persistance}
	\begin{itemize}
		\item Persistence unit définie dans \texttt{META-INF/\href{http://xmlns.jcp.org/xml/ns/persistence/persistence_2_1.xsd}{persistence.xml}} dans un jar comme bibliothèque du \texttt{.ear} {\tiny (aussi possible dans EJB ou .war)}
%		\item Cet EAR est alors le scope de l’unité
		\item Génération du schéma par JPA : dans \texttt{persistence.xml}, propriété \texttt{javax.persistence.schema-generation.database.action}, valeur \texttt{drop-and-create}
		\item Fournisseur trouvé dans classpath
%<property name="hibernate.show_sql" value="true"/>
%<property name="hibernate.format_sql" value="true"/>
	\end{itemize}
\end{frame}

\begin{frame}
	\frametitle{Unité de persistance en Java SE}
	\begin{itemize}
		\item Persistence.createEntityManagerFactory("helloworld");
	\end{itemize}
%<properties>
%<property name="hibernate.connection.driver_class"
%value="org.hsqldb.jdbcDriver"/>
%<property name="hibernate.connection.url"
%value="jdbc:hsqldb:hsql://localhost"/>
%<property name="hibernate.connection.username"
%value="sa"/>
\end{frame}

\begin{frame}
	\frametitle{Unité de persistance en Java EE}
	\begin{itemize}
		\item En général, préciser jta-data-source : un nom JNDI
		\item Conteneur Java EE définit \texttt{java:comp/DefaultDataSource} : nom utilisé par défaut
	\end{itemize}
\end{frame}

%instancier EM en Java SE : ben, l’instancier, et puis le fermer ! Démarrer et fermer la transaction entretemps.

%\subsection[Transactions EE]{Transactions et Java EE}
\begin{frame}
	\frametitle{Transactions et Java EE}
	\begin{itemize}
		\item Pilotes JDBC supportent {\tiny généralement} JTA (Java Transaction API)
		\item Permet {\tiny entre autres} la gestion des transactions par le conteneur
		\item Transactions gérées par le conteneur par défaut pour EJB {\tiny (ou utiliser \jeeref[@]{javax.ejb/TransactionManagement})}
		\item Chaque méthode participe alors par défaut à une transaction
		\item Si pas de transaction en cours lors de l’appel : le conteneur démarre puis termine automatiquement une transaction (commit si ok, rollback si exception {\tiny sauf exception application})
		\item Si transaction en cours, la méthode y participe
		\item Ou annoter la méthode \jeeref[@]{javax.ejb/TransactionAttribute}
	\end{itemize}
\end{frame}

\section{Exercices}
\begin{frame}
	\frametitle{Exercices}
	\begin{itemize}
		\item Créer une entité pour un type de votre projet.
		\item Faire en sorte que la table correspondante soit créée automatiquement lors du déploiement de l’entité.
		\item Permettre CR.D : Create, Retrieve, Delete \emph{aussi simple que possible}, via un ou plusieurs servlets. (N’utilisez pas de paramètres complexes, ce n’est pas le but de cet exercice.)
		\item Programmer une méthode qui transforme un attribut d’un objet.
		\item Permettre l’application de cette méthode via un SLSB appelé par un servlet. Votre servlet ne doit pas nécessairement accepter de paramètres. Quid de l’atomicité de la transaction ?
	\end{itemize}
\end{frame}

\section{JPA : Mise en œuvre}
\subsection{Liens entre entités}
\begin{frame}
	\frametitle{Liens entre entités}
	Exemple : un item a plusieurs bid associés {\tiny tiré de Java Persistence with Hibernate}
	\begin{itemize}
		\item Chaque bout d’une association : to one ou to many
		\item To one (\texttt{Item forItem} dans \texttt{Bid}) : annoter \jeeref[@]{javax.persistence/ManyToOne}, préciser éventuellement \texttt{optional=false} {\tiny et \texttt{fetch}}
		\item Si bidirectionnelle : votre code doit maintenir la synchronisation, malgré annotations (Pourquoi ? \pause pour indépendance à ORM)\pause
		\item Pour autre bout d’un lien bidirectionnel (ne compte pas dans DDL) : \jeeref[@]{javax.persistence/OneToMany} et préciser \texttt{mappedBy}
		\item Liens one to one similaire, utiliser \jeeref[@]{javax.persistence/OneToOne}
	\end{itemize}
\end{frame}

\subsection{Value type}
\begin{frame}
	\frametitle{Exemple pour Value type}
	\begin{itemize}
		\item User a une adresse
		\item Dans BD, une seule table
		\item Dans modèle objet, objet User et objet Adresse
		\item Modèle objet : granularité plus fine
		\item Pourquoi cette différence ? \pause
		\item Classes ont des responsabilités
		\item Réutiliser une classe, ne pas la dupliquer
		\item Schéma BD stable dans le temps
	\end{itemize}
\end{frame}

%\begin{frame}
%	\frametitle{}
%	\begin{itemize}
%		\item 
%	\end{itemize}
%	\begin{block}{}
%		
%	\end{block}
%\end{frame}
%
%\begin{frame}
%	\frametitle{}
%	\begin{itemize}
%		\item 
%	\end{itemize}
%	\begin{block}{}
%		
%	\end{block}
%\end{frame}
%
%\begin{frame}
%	\frametitle{}
%	\begin{itemize}
%		\item 
%	\end{itemize}
%	\begin{block}{}
%		
%	\end{block}
%\end{frame}

\subsection{Aspects pratiques}
\begin{frame}
	\frametitle{Concurrence}
	\begin{itemize}
		\item Entité avec \texttt{@Version} : protection optimiste par défaut
		\item Après update, effacement, merge : vérification automatique du n° de version en cache
		\item Vérification effectuée au moment même ou au \texttt{flush}
		\item Si versions ne correspondent pas : \jeeref{javax.persistence/OptimisticLockException}
		\item Protection pessimiste : utiliser \texttt{EntityManager.lock(entity, \href{https://docs.oracle.com/javaee/7/api/javax/persistence/LockModeType.html}{lockModeType})}
	\end{itemize}
\end{frame}

\begin{frame}
	\frametitle{Différentes égalités}
	\begin{block}{Trois types d’égalités}
		\begin{itemize}
			\item Égalité en mémoire : \texttt{a == b}
			\item Égalité objet : \texttt{a.equals(b)}
			\item Égalité DB : \texttt{a.getId().equals(b.getId())}
		\end{itemize}
	\end{block}
	Quelles propriétés prendre en compte dans \texttt{hashCode} et \texttt{equals} ?
	\begin{itemize}
		\item Table de hachage : objet ne peut changer de hash / d’égalité
		\item Dans une session il faut éviter deux objets not equals() et concernant la même ligne de la table
		\item Id pour égalité : pourquoi pas ? \pause Tant que non persistantes, ne fonctionne pas ; change lors sauvegarde \pause
		\item[$⇒$] Ensemble d’attributs déterminants pt de vue utilisateur ! (username…)
	\end{itemize}
\end{frame}

\begin{frame}
	\frametitle{Patterns}
	\begin{itemize}
		\item DM \emph{sans} dépendance persistance (tests ;  simplification dépendances)
		\item Entités transversales (couche web et business)
		\item Qqs classes service business
		\item Qqs classes persistence visibles du business
		\item Ou pattern DAO…
		\item Explorer : SFSB et \jeeref[@]{javax.persistence/PersistenceContext}(\jeeref{javax.persistence/PersistenceContextType\#EXTENDED})
	\end{itemize}
\end{frame}

\begin{frame}
	\frametitle{Usages plus avancés et divers}
	\begin{itemize}
		\item Entité peut être non chargée entièrement (\texttt{getReference})
		\item Accès éventuellement impossible après fermeture du contexte
		\item Contexte peut persister au-delà de la transaction ; peut être non synchronisé (\jeeref[@]{javax.persistence/PersistenceContext})
		\item Transitivité de la persistence, cf. \texttt{cascade} sur \texttt{@OneToMany} (par exemple)
		\item Type \og{}Value\fg{} plutôt que Entity pour des objets de cycle de vie dépendents d’autres
	\end{itemize}
	
\end{frame}

\subsection{Références}
\begin{frame}
	\frametitle{Références}
	\begin{itemize}
		\item Java Persistence with Hibernate : \href{http://gen.lib.rus.ec/book/index.php?md5=5D9F8BC8761804C0EBB8FE6A60BCF817}{1\iere} édition, \href{https://www.manning.com/books/java-persistence-with-hibernate-second-edition}{2\ieme} édition
		\item Hibernate 5.1 \href{http://docs.jboss.org/hibernate/orm/5.1/userguide/html_single/Hibernate_User_Guide.html}{User Guide}
		\item \href{https://jcp.org/en/jsr/detail?id=338}{JSR 338} (JPA 2.1) (\href{http://download.oracle.com/otn-pub/jcp/persistence-2_1-fr-eval-spec/JavaPersistence.pdf}{direct})
		\item \href{https://jcp.org/en/jsr/detail?id=907}{JSR 907} (JTA) (moins utile pour développeur Java EE)
		\item ENORM: An Essential Notation for Object-Relational Mapping, ACM SIGMOD Record 43(2), June 2014, pp 23–28 (\href{http://dx.doi.org/10.1145/2694413.2694418}{doi}, \href{http://www.sigmod.org/publications/sigmod-record/1406/pdfs/05.articles.Torres.pdf}{pdf} article)
%	Patterns Command, tout ça
	\end{itemize}
\end{frame}

\appendix
\AtBeginSection{
}
\section{Licence}
\begin{frame}
	\frametitle{Licence}
	Cette présentation, et le code LaTeX associé, sont sous \href{http://opensource.org/licenses/MIT}{licence MIT}. Vous êtes libres de réutiliser des éléments de cette présentation, sous réserve de citer l’auteur.
	
	Le travail réutilisé est à attribuer à \href{http://www.lamsade.dauphine.fr/~ocailloux/}{Olivier Cailloux}, Université Paris-Dauphine.
	
	\small{(Ceci ne couvre pas les images incluses dans ce document, puisque je n’en suis généralement pas l’auteur.)}
\end{frame}

\end{document}
\begin{frame}
	\frametitle{}
	\begin{itemize}
		\item 
	\end{itemize}
	\begin{block}{}
		
	\end{block}
\end{frame}

\section{Bibliographie}
\begin{frame}[allowframebreaks]
	\frametitle{Bibliographie}
	\def\newblock{\hskip .11em plus .33em minus .07em}
% 	\bibliography{zotero}
\end{frame}

\section{Autres}
\begin{frame}
	\frametitle{}
	\begin{itemize}
		\item 
	\end{itemize}
	\begin{block}{}
		
	\end{block}
\end{frame}
\end{document}
